\subsection{Observable and Hidden Causal Sectors}

The SBC framework assumes that the total causal structure can be decomposed into an observable sector and a hidden sector,
\begin{equation}
\Psi_{\rm tot}
=
\Psi_{\rm obs}
+
\Psi_{\rm hid}.
\end{equation}

The observable component is obtained through a projection operator,
\begin{equation}
\Psi_{\rm obs}
=
\mathcal P \Psi_{\rm tot},
\end{equation}
while the hidden component is
\begin{equation}
\Psi_{\rm hid}
=
(1-\mathcal P)\Psi_{\rm tot}.
\end{equation}

The projection operator satisfies
\begin{equation}
\mathcal P^2=\mathcal P,
\end{equation}
which defines a consistent separation between accessible and inaccessible causal sectors.

The total state is assumed to evolve according to
\begin{equation}
\frac{d\Psi_{\rm tot}}{dt}
=
\mathcal L \Psi_{\rm tot},
\end{equation}
where \(\mathcal L\) is the generator of the underlying causal evolution.

Projecting the total dynamics onto the observable and hidden sectors yields
\begin{equation}
\frac{d\Psi_{\rm obs}}{dt}
=
\mathcal L_{oo}\Psi_{\rm obs}
+
\mathcal L_{oh}\Psi_{\rm hid},
\end{equation}
and
\begin{equation}
\frac{d\Psi_{\rm hid}}{dt}
=
\mathcal L_{ho}\Psi_{\rm obs}
+
\mathcal L_{hh}\Psi_{\rm hid}.
\end{equation}

Here,
\begin{equation}
\mathcal L_{oo}
=
\mathcal P\mathcal L\mathcal P,
\end{equation}
\begin{equation}
\mathcal L_{oh}
=
\mathcal P\mathcal L(1-\mathcal P),
\end{equation}
\begin{equation}
\mathcal L_{ho}
=
(1-\mathcal P)\mathcal L\mathcal P,
\end{equation}
and
\begin{equation}
\mathcal L_{hh}
=
(1-\mathcal P)\mathcal L(1-\mathcal P).
\end{equation}

The hidden sector acts as a reservoir of causal information that is not directly resolved within the observable window.

Solving formally for \(\Psi_{\rm hid}\) and substituting it back into the observable-sector equation generates an effective nonlocal equation,
\begin{equation}
\frac{d\Psi_{\rm obs}}{dt}
=
\mathcal L_{\rm eff}\Psi_{\rm obs}
+
\int_{-\infty}^{t}
K(t,t')
\Psi_{\rm obs}(t')
\,dt'.
\end{equation}

The memory kernel \(K(t,t')\) is therefore not introduced as an independent postulate. It emerges from eliminating the hidden causal sector.

This gives the fundamental SBC hierarchy:
\begin{equation}
{\rm total\ causal\ structure}
\rightarrow
{\rm causal\ projection}
\rightarrow
{\rm hidden\ sector}
\rightarrow
{\rm memory\ kernel}
\rightarrow
{\rm infrared\ phenomenology}.
\end{equation}

The SBC/F3 sector can therefore be interpreted as the cosmological infrared limit of an observable-hidden causal decomposition.
